\chapter{到達点}

\section{何が確立され、何がされなかったか}

\subsection{理論として導いたもの}

\begin{center}
\begin{tabularx}{\textwidth}{@{}lXl@{}}
\toprule
主張 & 内容 & 状態 \\
\midrule
$\Phi$ による記述 & 履行と決済の写像として書ける & 有効に機能した \\
$\kappa$ の統合性 & 売掛金・前受金・買掛金が単一の量になる & 有効 \\
$\kappa_i$ による統合 & 四分野の理論が同一の和に入る & 本稿の中心的な寄与（注\ref{rem:disjoint-literatures}） \\
$g^\star = (m+d-I/r)/\CCC$ & 自己金融可能成長率 & 導出済（系\ref{cor:gstar}） \\
$\CCC=\sum\tau_i+\tau_{\mathrm{inv}}$ & $\CCC$ は遅れ時間の和 & 導出済（命題\ref{prop:ccc-structure}） \\
容量制約下の会員制 & 低稼働が実行可能性の条件 & 導出済（命題\ref{prop:capacity}） \\
$E$ と $\phi$ の接続 & 内部留保は何の累積か & 定義\ref{def:equity}により $\phi$ の累積 \\
$\kappa_K$ の主体依存性 & 出資は他の $\kappa$ と同じか & 測度に依存し恒等式に入らない（命題\ref{prop:kappa-K-excluded}） \\
\bottomrule
\end{tabularx}
\captionof{table}{理論として導出した主張とその状態。}
\label{tab:theory-achievements}
\end{center}

\subsection{実証で判定したもの}

\begin{center}
\begin{tabularx}{\textwidth}{@{}lXl@{}}
\toprule
主張 & 内容 & 状態 \\
\midrule
$\CCC$ の符号 & 業種・媒体で二桁の水準差 & 実証で確認（第\ref{part:results}部） \\
$g^\star$ の推移 & 25年間の変化 & 低下していない（命題\ref{prop:gstar-flat}） \\
$g^\star$ の分子 & $m$ のみか投資も要るか & $I$ を含めると水準が大きく動く \\
小規模層の $m$ & 規模別の利益率格差 & 役員報酬で説明される（命題\ref{prop:officer}） \\
個人企業との比較 & 会計区分か経済実態か & 別統計で命題\ref{prop:officer}を支持（第\ref{sec:individual-firms}節） \\
$\CCC$ 上昇の要因 & 構成変化か業種内変化か & 業種内が83\%。手形の消滅が主因 \\
分母効果の寄与 & $\CCC=W/r$ の機械的変動 & 17\%程度。支配的ではない \\
外部化の対価 & 手数料は吸収した機能に対応 & 部分的に支持（第\ref{sec:platform-fee}章） \\
制度層の作用形態 & 上限賦課か競争条件の操作か & 後者が近年の主流（注\ref{rem:institution-shift}） \\
\midrule
$\CCC$ の予測力 & 景気先行指標としての機能 & \textbf{棄却}（第\ref{part:industry}章） \\
規模と与信の向き & 小規模ほど与信の出し手 & \textbf{棄却、符号が逆} \\
$\phi_{\rm cog}$ の測定 & 認知余剰の実測 & 本稿では\textbf{未達}（第\ref{sec:phicog-measurable}節） \\
可測性と契約形態 & 観測可能性が $\Phi$ を決める & \textbf{未検証}（第\ref{sec:predB-illposed}節） \\
手数料率と離職率 & 成果のリスクが価格に反映されるか & \textbf{棄却}。307社で無相関（第\ref{sec:predX-result}節） \\
\bottomrule
\end{tabularx}
\captionof{table}{実証で判定した主張とその状態。棄却・未検証を含む。}
\label{tab:empirical-judgments}
\end{center}

記述の道具としての枠組みは機能した。予測の道具としては機能しなかった。
この二つは別の性質であり、本稿は前者のみを主張する。

\section{事前登録が働いた記録}

本稿では16の予測を事前に文章として固定した。
支持されたのは部分的支持を含めて6件、棄却が6件、検証不能が4件である。

\begin{center}
\begin{tabularx}{\textwidth}{@{}llXl@{}}
\toprule
予測 & 対象 & 内容と結果 & 判定 \\
\midrule
A & 族2 & 乖離と取引頻度。水準は支持だが測れたのは $\kappa$ & 部分的支持 \\
B & 族5 & 返戻金と可測性。符号が一意に定まらない & 検証不能 \\
C & 手数料 & 外部化の対価。機能で説明できない反例も発見 & 部分的支持 \\
D & 無人化 & 依存の数と無人期間。方向は支持だが機構が誤り & 部分的支持 \\
E & 法人統計 & $\CCC$ の先行性（$p=0.72$） & 棄却 \\
F & 法人統計 & 規模と与信の向き。符号が逆 & 棄却 \\
J & 手数料 & 料率差は相関 $\rho$ に対応。$\rho$ が観測できない & 検証不能 \\
K & $\phi_{\rm cog}$ & $\delta^\ast>\delta$。四カテゴリで符号が負 & 棄却 \\
L & 容量制約 & 高齢層の契約選択。消費者側データなし & 検証不能 \\
M & 法人統計 & $\CCC$ 上昇の構成効果。業種内が83\% & 棄却 \\
N & 法人統計 & $\Delta\ln W$ と $\Delta\ln r$ の正相関（$\beta=0.83$） & 支持 \\
O & 法人統計 & 遅れ構造。年次データで検出できず & 棄却 \\
P & 法人統計 & ラグを含む総弾力性。$0.830\to0.821$ & 棄却 \\
T & 労働側 & 新法による $\kappa$ の移動。統計が存在しない & 検証不能 \\
U & 個人企業 & 営業利益率が役員報酬を戻した値に近い & 支持 \\
X & 族5 & 手数料率と離職率の負相関。307社で無相関 & 棄却 \\
\bottomrule
\end{tabularx}
\captionof{table}{事前登録した16予測の対象・内容・判定結果。}
\label{tab:predictions-16}
\end{center}

Q・R・S は労働者視点の検討過程で登録したが、
第\ref{sec:worker-predictions}節のとおり検証に着手していない。
G・H・I は制度による料率規制を対象としたが、
母集団が3件しかないことが判明したため設計段階で取り下げた。

事前登録が具体的に働いた箇所を記録しておく。

\begin{itemize}
\item 族5において、上限制手数料が99.7\%消滅している事実は、
可測性仮説の傍証として解釈することが可能であった。
実際には上限制の価格水準が一桁低いだけで無関係である。
予測を事前に固定していなければ、後付けで理論の傍証に採用していた可能性が高い。

\item 第\ref{part:industry}章において、$\CCC$ と売上高の同期負相関を
「$\CCC$ は景気指標として機能する」と読み、
大企業の与信ポジションを「大企業が取引先を支えている」と解釈しえた。
いずれも後付けの物語である。
\end{itemize}

第\ref{sec:narrative}節で他者の記述について述べた事後的物語化は、
自らの解釈にも生じる。

\section{本稿の途中で行った訂正}

検討の過程で撤回した主張が12件ある。内訳と契機は付録\ref{app:revisions}に記録した。
最も重いのは $g^\star$ に関する訂正である。
$\CCC$ の悪化のみを観測して自己金融可能成長率の低下を結論しかけたが、
系\ref{cor:gstar}は比であり、粗利率 $m$ を取得した結果、
$g^\star$ は低下していなかった（命題\ref{prop:gstar-flat}）。

$\phi_{\rm cog}$ と可測性仮説について「測定されていない」「予測が不良設定である」と
記していた二件は、先行研究の確認により不正確であることが判明した。
前者には測定例が存在し、後者は分野全体の未解決問題である。

\section{未解決の問題}

未解決の問題を、第\ref{sec:obstacle-taxonomy}節の分類に従って整理する。
本稿の調査により四分類に収まらない障害が二つ現れたため、それらも含める
（第\ref{sec:sixth-obstacle}節）。
分類が異なれば必要な対処も異なるため、労力の配分はこの区別に依存する。

\subsection{（i）測定の不在に属するもの}

当該の量が世界のどこでも測定されていない。二次データでは埋まらず、原調査を要する。

\begin{enumerate}[label=(\arabic*)]
\item \textbf{役員報酬は費用か利益処分か。}
小規模層の $m$ が低い原因は役員報酬である（命題\ref{prop:officer}）。
売上高比で大企業の50倍にあたり、営業利益に戻すと規模による序列が消失する。

残る問いは、この報酬が労働の対価か裁量的な利益処分かである。
利益が変動しても役員報酬率が動かないこと、
および個人企業の営業利益率が不況期に上昇すること（注\ref{rem:covid-reversal}）は
裁量性を示唆するが、労働対価説も棄却できない。
判定には役員の労働時間が必要だが、
\textbf{役員は労働基準法上の労働者ではないため、労働時間の公的統計が存在しない}
（注\ref{rem:officer-undecidable}）。

代替経路として法人税率の変更や軽減税率の適用区分を外生ショックとする識別が考えられるが、
閾値で標本を切るには個票を要する。
\end{enumerate}

\subsection{（ii）仮説の不良設定に属するもの}

理論が符号を一意に定めない。データを増やしても解決せず、理論の精密化を要する。

\begin{enumerate}[label=(\arabic*),start=2]
\item \textbf{可測性仮説の検定。}
$b^\star$ の式そのものは、必要な量が揃わないため検定できない
（第\ref{sec:what-to-measure}節）。

価格側からの検定は実施した。307社について
手数料実績率と離職率の相関を測ったが、無相関であった（第\ref{sec:predX-result}節）。
職種平均で離職率が5.5倍異なるのに手数料率は1.16倍しか異ならず、
価格は 20--30\% の慣行値に集中している。

したがって障害は（ii）仮説の不良設定だけではない。
\textbf{説明変数が理論の想定する範囲で変動していない}という
第\ref{sec:sixth-obstacle}節の（vi）に該当し、
理論の精密化でもデータの追加でも解決しない。

\item \textbf{R2 の目的関数：$H$ と $\beta$ で足りるか。}
目的関数の形を変えずに $H$（時間展望）と $\beta$（割引因子）の変化として
説明する経路を示した（第\ref{sec:r2-sst}節）。これが十分かは未検証である。

検証には予測L の検定が必要だが、年齢別の契約形態選択分布が得られていない。
供給側の価格設計は対抗仮説と整合するが、消費者の選択ではない。
\end{enumerate}

\subsection{（iii）識別不能に属するもの}

データはあるが交絡を分離できない。設計の工夫か外生変動を要する。

\begin{enumerate}[label=(\arabic*),start=4]
\item \textbf{売掛金伸長の解釈。}
$\CCC$ の長期上昇について、第\ref{sec:ccc-decomposition}節で要因を分解した。
産業構成の変化は排除され（構成効果は17\%、2010年以降は2\%）、
2000年代の上昇は手形の消滅による DPO 短縮で説明される。

残るのは2010年代の売掛金伸長（$+6.62$ 日）である。
実質的な回収期間の長期化か、電子記録債権への移行に伴う科目移動かを分離できない。
分離には電子記録債権の残高を要するが、法人企業統計の調査項目にない（注\ref{rem:densai}）。

\item \textbf{$\Delta\ln W$ の変動要因。}
同期負相関の分離は第\ref{sec:denominator-effect}節で行い、
分母効果の寄与は17\%程度と得た。分母効果が支配的でないことは示せた。

遅れ構造を考慮した推定も試みたが、年次データでは検出できなかった（注\ref{rem:lag-undetectable}）。
弾力性推定の $R^2$ は 0.081 にとどまり、
運転資本の年次変動の大半が売上以外の要因で駆動されている。
決算期の在庫調整、大口取引の期ずれ、会計方針の変更が混入すると考えられるが、分離していない。
\end{enumerate}

\subsection{（iv）アクセス制約に属するもの}

存在するが取得できない。経路の開拓か断念を要する。

\begin{enumerate}[label=(\arabic*),start=6]
\item \textbf{$\phi_{\rm cog}$ の測定。}
命題\ref{prop:phicog-requirements}の三条件のうち、
本稿が到達できたのは(1)契約上の権利のみである。
(2)個票の実現利用と(3)注意を強制する外生イベントを欠く。

迂回路として $\delta^{*}=c/p$ を構成したが、二択メニューにのみ適用でき、
階段状の料金体系では代理が一意に定まらない（注\ref{rem:detour-limit}）。

族2-4（プリペイド）については \cite{nber2026pnbl} が
前払い価値の約40\%が未使用であることを示しており、
\textbf{本稿が族2で測れなかった量は既に測定されている}。
残る空白は族2-5（オプション）についての推定と、
$\phi$ 全体に占める $\phi_{\rm cog}$ の比率の業種横断推計である。

\item \textbf{余剰の取りやすさの指標。}
第\ref{part:industry}章の二軸のうち横軸に対応する指標を構成できなかった。
断念の理由は構造的である。

\textbf{仲介者は業種を跨ぐ主体であり、業種の属性ではない}
（注\ref{rem:no-horizontal}）。業種別集計では原理的に捕捉できない。
仲介者側の断面を取ろうとすると、料率が個別審査で決まるため公開されない。
\textbf{料率が交渉で決まるとは、当事者以外に知られないことが前提である。}
公開料率のみを集めれば、公表料金表で運用する薄い仲介者に偏る。

リスク成分については予測J（第\ref{sec:pred-J}節）を立てたが、
$\rho$ の推定に要するデータが仲介者の内部にあり、
かつ $\rho$ の高い業種は審査で排除されるため検証できない。
\end{enumerate}

\subsection{枠組みの記述範囲に属するもの}

障害の分類に収まらない、本稿の対象設定そのものに関わる項目。

\begin{enumerate}[label=(\arabic*),start=8]
\item \textbf{容量制約の簡略化。}
注\ref{rem:cap-simplification}のとおり、本稿は $\mathrm{Cap}$ を定数として扱った。
第\ref{sec:cap-observation}節にこれが合わない観測を記録したが、
$n=1$ であるため理論は修正していない（注\ref{rem:no-generalization}）。
他の族でも同じ限界が現れるかは未確認である。

\item \textbf{集約の社会的最適性。}
本稿は個別事業者の視点で書かれており、
$\Phi$ の選択が系に与える影響を扱わない（注\ref{rem:pooling-social}）。
族5-4（保険引受）と族6-2（エスクロー）は集約を業とする類型であり、
個別に最適な集約が系には過剰でありうる。

\item \textbf{個人事業主の信用ポジション。}
個人事業主の $\kappa$ を測る公表統計が存在しない。
個人企業経済調査は2019年度に資産・負債の調査を廃止し、
中小企業実態基本調査は個人事業者用の調査票を持つが資産・負債は集計されない
（第\ref{sec:not-tabulated}節）。

これにより、フリーランス新法（2024年11月施行）が
$\tau$ の上限を規定したことの効果を検証できない。
第\ref{sec:institution-mode}節で整理した制度層の三形式に対し、
これは\textbf{時間を直接規定する第四の形式}であり、
第\ref{sec:tau}節の $\tau_i$ に法が上限を課す事例として理論的な関心が高い。

\item \textbf{請求権が流通する形式。}
第\ref{sec:phi}章の $\Phi:\Delta\to\Pi$ は、
履行と決済が同一の相手方との間で対応することを前提とする。
トークン発行のように決済が第三者へ転売されると
履行を受ける主体が事後的に変わり、$N$ を固定した設定では表現できない
（第\ref{sec:token-limit}節）。定義域の拡張を要する。

なお当初この項目に含めていた労働側の $\Phi$ と B2B 複合契約は、
既存類型で記述できることが確認された（第\ref{sec:labor-phi}節）。
\end{enumerate}

\subsection{優先順位}

実行可能性、波及範囲、費用の三点で評価すると、次の順になる。

\begin{center}
\begin{tabularx}{\textwidth}{@{}rXlll@{}}
\toprule
 & 項目 & 分類 & 実行可能性 & 波及 \\
\midrule
1 & (3) R2 の $H$ と $\beta$ & (ii) & 中 & 中 \\
2 & (5) $\Delta\ln W$ の変動要因 & (iii) & 中 & 中 \\
3 & (10) 個人事業主の $\kappa$ & 集計の選択 & 中 & 大 \\
4 & (4) 売掛金伸長の解釈 & (iii) & 低 & 小 \\
5 & (6) $\phi_{\rm cog}$ & (iv) & 低 & 大 \\
6 & (1) 役員報酬 & (i) & 低 & 大 \\
7 & (2) 可測性仮説 & (ii) & 低 & 中 \\
8 & (7) 余剰の取りやすさ & (iv) & 低 & 中 \\
9 & (9) 集約の社会的最適性 & --- & 低 & 中 \\
10 & (8) 容量制約 & --- & 低 & 小 \\
11 & (11) 請求権の流通 & --- & 低 & 中 \\
\bottomrule
\end{tabularx}
\captionof{table}{未解決問題11項目の分類・実行可能性・波及範囲による優先順位。}
\label{tab:unresolved-priority}
\end{center}

（i）と（ii）に属する項目は労力の投入では前進しない。
実行可能性が中位にあるのは(3)(5)(10)の三項目である。

このうち(10)個人事業主の $\kappa$ は性質が異なる。
障害は測定の不在でも識別不能でもなく\textbf{集計の選択}であり、
既に収集されている可能性のある回答が公表されないことによる。
第\ref{sec:not-tabulated}節のとおり中小企業実態基本調査には
個人事業者用の調査票が存在する。
統計作成者に対する集計要望、または匿名データの利用申請という経路がありうる。

\textbf{公表データの取得のみで前進できる項目は、現時点で残っていない。}

\section{枠組みの既知の弱点}

検証を通じて明らかになった弱点は、三つの類型に整理できる。
いずれも個別の命題の誤りではなく、\textbf{枠組みの構成に由来する}。

\subsection{（A）独立性の過剰仮定}
\label{sec:independence-assumption}

複数の命題が、実際には同時に決まる量を独立として扱っている。

\begin{center}
\begin{tabularx}{\textwidth}{@{}lXX@{}}
\toprule
命題 & 独立と扱った組 & 実際の関係 \\
\midrule
\ref{prop:pooling} 集約 & $n$ と $\rho$ & 規模拡大は相関の低い対象を枯渇させる（注\ref{rem:rho-endogenous}） \\
\ref{prop:breakage} 乖離 & $\bar\delta$ と $\delta$ & メニュー選択を通じた同時決定 \\
\ref{prop:capacity} 会員制 & $m$ と $\E[\delta]$ & 会費が両方を動かす \\
\ref{prop:implementable} 実装可能性 & $\mathcal{G}$ と $\pi$ & $\mathcal{G}$ は投資対象（注\ref{rem:G-endogenous}） \\
\ref{prop:fcf} FCF & $m$ と $\CCC$ & 実測で同方向に上昇（注\ref{rem:m-ccc-independence}） \\
\ref{prop:substitution} 代替性 & $E$ と $\kappa$ & 資本があれば与信を出せる（命題\ref{prop:officer}の系） \\
\bottomrule
\end{tabularx}
\captionof{table}{独立性を仮定した命題と、実際には同時に決まる量との対応。}
\label{tab:independence-assumptions}
\end{center}

各命題には独立性の仮定を明示したが、\textbf{仮定を外した場合の帰結は導いていない}。
証明は形式的には正しく、適用範囲が狭いという性質の弱点である。

事前登録した予測のうち四件が、この弱点に由来して外れた。

\begin{center}
\begin{tabularx}{\textwidth}{@{}lXX@{}}
\toprule
予測 & 独立と扱った組 & 誤りの形 \\
\midrule
A & $\kappa$ と $\phi_{\rm cog}$ & 滞留と非行使を同一視した \\
B & $b$ と $\sigma^2$ & 関係式の分子側のみを見た \\
F & DSO と DPO & 交渉力が両側に効くことを見落とした \\
K & $\bar\delta$ と $\delta$ & メニューの階段構造を想定しなかった \\
\bottomrule
\end{tabularx}
\captionof{table}{独立性の過剰仮定に由来して外れた4件の予測。}
\label{tab:independence-failed-predictions}
\end{center}

BとFは同型である。\textbf{両側に効く変数を片側だけで扱っている}。

\subsection{（B）外生性の過剰仮定}

制度、$\Phi$、$v$、$\mathrm{Cap}$、$\mathcal{G}$ を外生として扱っている。

\begin{enumerate}[label=(\arabic*)]
\item \textbf{評価写像 $v$ の主観性。}
履行を貨幣で測る $v$ を所与としたが、これ自体が主観的であり、
$\phi_{\mathrm{cog}}$ の定義がこれに依存している。

\item \textbf{$\Phi$ の内生性と切り替え費用。}
$\Phi$ を企業の選択変数として扱ったが、
実際には主体間の交渉の均衡解であり、単独では選択できない。
加えて切り替え費用をゼロと置いている（注\ref{rem:switching-cost}）。
費用が余剰の増分と同程度であれば待機の価値が生じ、
$\Phi$ をいつ確定させるかというタイミングの問題が現れる。
本稿はこれを扱わない。

\item \textbf{制度層の外生性。}
制度を外生的制約としてのみ扱った。
実際には制度は企業行動に反応して内生的に変化する。
第\ref{sec:platform-fee}章の反トラスト訴訟と料率引き下げの対応がその例である。

\item \textbf{$\mathrm{Cap}$ の簡略化。}
定数として扱ったが、本来は時刻の関数であり、
かつ利用時間帯の制限を通じて設計変数でもある（注\ref{rem:cap-simplification}）。

\item \textbf{$\mathcal{G}$ の内生性。}
情報構造を所与としたが、事業者は追跡調査により拡張できる（注\ref{rem:G-endogenous}）。
\end{enumerate}

\subsection{（C）水準と変化の未分離}
\label{sec:level-change-gap}

\textbf{最も見落とされやすい弱点である。}
(A) と (B) は「仮定が強すぎる」という形で認識できるが、
(C) は理論部に対応する概念が存在しないため、欠落として気づきにくい。

注\ref{rem:level-vs-change}のとおり、本稿の量はすべて水準として定義されている。
ところが実証では、水準の横断比較が解釈できず、変化でしか語れない事態に四度直面した。

\begin{center}
\begin{tabularx}{\textwidth}{@{}lX@{}}
\toprule
局面 & 対処 \\
\midrule
族2 媒体別パネル & 水準差は二桁。成長セグメントの併存で時点効果を分離 \\
第\ref{part:industry}章 業種別 $\CCC$ & 5--210日。横断比較は解釈不能と明記 \\
予測M シフト・シェア分解 & 水準変化を業種内効果と構成効果に分解 \\
予測N 弾力性推定 & 変化同士の関係として $\beta$ を推定 \\
\bottomrule
\end{tabularx}
\captionof{table}{水準と変化の未分離が現れた局面と、それぞれの対処。}
\label{tab:level-change-gap-response}
\end{center}

\textbf{理論部が水準しか持たないのに、実証は変化でしか語れない。}

この不整合の一部は、第\ref{sec:tau}節で遅れ時間 $\tau_i$ を導入することで解消した。
命題\ref{prop:ccc-structure}により $\CCC=\sum_i\tau_i+\tau_{\mathrm{inv}}$ となり、
$\CCC$ の水準が $\Phi$ に決まること（系\ref{cor:ccc-level}）と、
横断比較が無意味であること（系\ref{cor:no-cross-comparison}）が演繹される。
$\Delta\CCC$ を契約の変化と運用の変化に分ける道具も得られた（系\ref{cor:ccc-change}）。

\textbf{残るのは比較可能な量の構成である。}
候補として契約上の $\tau_i^{\mathrm{contract}}$ からの乖離が挙がるが、
契約条件は公表統計に存在しない（注\ref{rem:no-normalization}）。
理論の問題は解けたが、測定の問題に移った。

\subsection{三類型と中粒度理論}

(A) と (C) は、接続すべき既存理論が特定できる。

\begin{center}
\begin{tabularx}{\textwidth}{@{}lXX@{}}
\toprule
弱点 & 該当箇所 & 接続しうる理論 \\
\midrule
(A) $\bar\delta$ と $\delta$ & 命題\ref{prop:breakage}、族2、フィットネス & メカニズムデザイン、自己選択と数量割引 \\
(A) $b$ と $\sigma^2$ & 命題\ref{prop:implementable}、族5 & 契約の選別効果 \\
(A) DSO と DPO & 命題\ref{prop:substitution}、第\ref{part:industry}章 & 取引信用（trade credit）の理論 \\
(C) 水準と変化 & 全実証章 & --- \\
\bottomrule
\end{tabularx}
\captionof{table}{枠組みの弱点と、接続しうる既存理論の対応。}
\label{tab:weakness-mid-level-theory}
\end{center}

\textbf{本稿はこれらを参照したが、体系的に接続していない。}
第\ref{sec:info}章で情報性原理を借りながら、
第\ref{sec:predB-illposed}節で予測を立てる際にその論争状況を確認しなかったのは、
同じ理論を参照するときと予測を立てるときで異なる深さで扱っていたためである。

(C) については、接続先を特定できていない。

